\documentclass[12pt, a4paper]{article}
% --- 1. Cấu hình Tiếng Việt và Font chữ chuẩn ---
\usepackage[utf8]{inputenc}
\usepackage[T5]{fontenc}
\usepackage[vietnamese]{babel}
\usepackage{mathptmx} % Font Times New Roman đồng bộ cả chữ và công thức toán, không gây xung đột gói

\makeatletter
\renewcommand{\normalsize}{\@setfontsize\normalsize{13pt}{16pt}}
\makeatother
\normalsize

% --- 2. Gói Toán học & Ký hiệu bổ trợ ---
\usepackage{amsmath, amssymb, amsfonts, mathrsfs, amsthm}
\usepackage{graphicx}
\usepackage{fontawesome}
\usepackage{booktabs, tabularx, array, multirow}
\usepackage{enumitem}

% --- 3. Đồ họa TikZ, Bảng biến thiên & Hình không gian ---
\usepackage{tikz}
\usetikzlibrary{calc, intersections, angles, quotes, patterns, arrows.meta, 3d}
\usepackage{pgfplots}
\pgfplotsset{compat=1.18}
\usepackage{tkz-euclide}
\usepackage{tkz-tab}

\renewcommand{\vec}[1]{\overrightarrow{#1}}

% --- 4. Hộp sư phạm (Tcolorbox) ---
\usepackage[many]{tcolorbox}
\tcbuselibrary{skins, breakable}

% --- 5. Căn lề chuẩn văn bản giáo án ---
\usepackage{geometry}
\geometry{
    a4paper,
    left=25mm,
    right=15mm,
    top=20mm,
    bottom=20mm
}

% --- 6. Header / Footer chuẩn hóa ---
\usepackage{fancyhdr}
\usepackage{lastpage}
\pagestyle{fancy}
\fancyhf{}

\fancyhead[L]{\small \textit{Kế hoạch bài dạy Toán 11}}
\fancyhead[R]{\textbf{Trang \thepage/\pageref{LastPage}}}
\fancyfoot[L]{\small \textit{Tổ Toán - Tin}}
\fancyfoot[R]{\small \textit{Trường THCS\&THPT Hồng Vân}}
\renewcommand{\headrulewidth}{0.4pt}
\renewcommand{\footrulewidth}{0.4pt}

% --- 7. Giãn dòng & Giãn đoạn theo chuẩn Nghị định 30/2020/NĐ-CP ---
\usepackage{setspace}
\setstretch{1.15}
\setlength{\parindent}{1.0cm}
\setlength{\parskip}{5pt}

% Định nghĩa hộp sư phạm chuyên dụng
\newtcolorbox{pedabox}[2][]{
    enhanced,
    breakable,
    colback=blue!3!white,
    colframe=blue!65!black,
    fonttitle=\bfseries\small,
    title={#2},
    arc=2mm,
    boxrule=0.6pt,
    left=3mm, right=3mm, top=2mm, bottom=2mm,
    #1
}

\newtcolorbox{aibox}[2][]{
    enhanced,
    breakable,
    colback=teal!4!white,
    colframe=teal!70!black,
    fonttitle=\bfseries\small,
    title={#2},
    arc=2mm,
    boxrule=0.6pt,
    left=3mm, right=3mm, top=2mm, bottom=2mm,
    #1
}

\newtcolorbox{scaffoldbox}[2][]{
    enhanced,
    breakable,
    colback=orange!4!white,
    colframe=orange!80!black,
    fonttitle=\bfseries\small,
    title={#2},
    arc=2mm,
    boxrule=0.6pt,
    left=3mm, right=3mm, top=2mm, bottom=2mm,
    #1
}

\begin{document}

\begin{center}
    {\fontsize{14pt}{17pt}\selectfont \textbf{KẾ HOẠCH BÀI DẠY MÔN TOÁN LỚP 11}}\\[6pt]
    {\fontsize{16pt}{19pt}\selectfont \textbf{BÀI 1. GIÁ TRỊ LƯỢNG GIÁC CỦA GÓC LƯỢNG GIÁC}}\\[4pt]
    {\fontsize{12pt}{15pt}\selectfont \textbf{Môn học: Toán 11} \ \textbar \ \textbf{Thời lượng: 3 tiết (Tiết 1 đến Tiết 3 theo PPCT)}}
\end{center}

\vspace{5pt}

\section*{I. MỤC TIÊU}

\subsection*{1. Kiến thức, Năng lực}

\begin{itemize}[leftmargin=1.2cm]
    \item \textbf{Yêu cầu cần đạt (Nguyên văn CT GDPT 2018):}
    \begin{itemize}[leftmargin=0.6cm]
        \item Nhận biết được các khái niệm cơ bản về góc lượng giác: khái niệm góc lượng giác; số đo của góc lượng giác; hệ thức Chasles cho các góc lượng giác; đường tròn lượng giác.
        \item Nhận biết được khái niệm giá trị lượng giác của một góc lượng giác.
        \item Mô tả được bảng giá trị lượng giác của một số góc lượng giác thường gặp; hệ thức cơ bản giữa các giá trị lượng giác của một góc lượng giác; quan hệ giữa các giá trị lượng giác của các góc lượng giác có liên quan đặc biệt: bù nhau, phụ nhau, đối nhau, hơn kém nhau $\pi$.
        \item Sử dụng được máy tính cầm tay để tính giá trị lượng giác của một góc lượng giác khi biết số đo của góc đó.
    \end{itemize}

    \item \textbf{Năng lực Toán học đặc thù:}
    \begin{itemize}[leftmargin=0.6cm]
        \item \textit{Năng lực tư duy và lập luận toán học:} So sánh được sự khác biệt bản chất giữa góc hình học (giới hạn từ $0^\circ$ đến $180^\circ$) và góc lượng giác (chiều quay, số đo tùy ý); phân tích dấu hoành độ, tung độ của điểm biểu diễn trên đường tròn lượng giác để suy ra dấu của các giá trị lượng giác $\sin, \cos, \tan, \cot$ trên bốn góc phần tư; lập luận chứng minh các hệ thức lượng giác cơ bản.
        \item \textit{Năng lực mô hình hóa toán học:} Thiết lập mô hình góc lượng giác và đường tròn lượng giác để biểu diễn chuyển động quay của kim đồng hồ, đu quay Mặt Trời, chuyển động lăn của bánh xe đạp.
        \item \textit{Năng lực giải quyết vấn đề toán học:} Vận dụng các hệ thức lượng giác cơ bản và công thức góc liên quan đặc biệt để tính giá trị lượng giác còn lại khi biết một giá trị lượng giác, rút gọn biểu thức, chứng minh đẳng thức.
        \item \textit{Năng lực giao tiếp toán học:} Trình bày rõ ràng khái niệm tia đầu, tia cuối, chiều quay, điểm biểu diễn của góc lượng giác; giải thích logic các bước biến đổi công thức.
        \item \textit{Năng lực sử dụng công cụ, phương tiện học toán:} Sử dụng thước kẻ, compa vẽ đường tròn lượng giác; sử dụng thành thạo máy tính cầm tay (MTCT) Casio fx-580VN X để đổi đơn vị góc và tính các giá trị lượng giác.
    \end{itemize}

    \item \textbf{Năng lực số (Theo Thông tư 02/2025/TT-BGDĐT):}
    \begin{itemize}[leftmargin=0.6cm]
        \item \textit{Mã 5.2NC1b:} Học sinh áp dụng công cụ máy tính cầm tay (MTCT Casio fx-580VN X) thành thạo chuyển đổi giữa đơn vị độ và radian (cài đặt \texttt{SHIFT MENU 2 1} / \texttt{SHIFT MENU 2 2}), tính toán chính xác giá trị lượng giác của các góc bất kỳ.
        \item \textit{Khai thác và tương tác phần mềm số:} Học sinh quan sát, tương tác với mô phỏng động GeoGebra về đường tròn lượng giác để kiểm chứng trực quan mối liên hệ giữa vị trí điểm $M$ và giá trị $\cos\alpha, \sin\alpha$.
    \end{itemize}

    \item \textbf{Năng lực AI (Theo Quyết định 2422/QĐ-BGDĐT):}
    \begin{itemize}[leftmargin=0.6cm]
        \item \textit{Mã 11.C3.1:} Học sinh xác định được một số kĩ thuật prompt đặt câu lệnh hỏi trợ lý AI giải thích trực quan về đường tròn lượng giác, kiểm chứng các giá trị lượng giác đặc biệt.
        \item \textit{Đánh giá và nhận diện giới hạn của AI:} Học sinh nhận thức được giới hạn của AI khi nhận diện đơn vị góc (dễ nhầm lẫn giữa độ và radian nếu prompt không nêu rõ ngữ cảnh) và biết cách đối chiếu kết quả của AI với máy tính cầm tay và bảng lượng giác chuẩn.
    \end{itemize}

    \item \textbf{Năng lực chung:}
    \begin{itemize}[leftmargin=0.6cm]
        \item \textit{Tự chủ và tự học:} Tự giác nghiên cứu tài liệu, hoàn thành các nhiệm vụ trên phiếu học tập cá nhân.
        \item \textit{Giao tiếp và hợp tác:} Tương tác tích cực trong thảo luận cặp đôi và nhóm nhỏ, phân công nhiệm vụ rõ ràng.
    \end{itemize}
\end{itemize}

\subsection*{2. Phẩm chất}

\begin{itemize}[leftmargin=1.2cm]
    \item \textbf{Chăm chỉ:} Kiên trì thực hiện đầy đủ các phép tính toán, vẽ hình chính xác, tích cực luyện tập chuyển đổi đơn vị.
    \item \textbf{Trung thực:} Khách quan, trung thực trong quá trình tự đánh giá, báo cáo kết quả bài làm cá nhân và nhóm.
    \item \textbf{Trách nhiệm:} Có trách nhiệm hỗ trợ bạn trong nhóm (đặc biệt học sinh trung bình -- yếu), bảo quản cẩn thận phương tiện học tập.
\end{itemize}

\section*{II. THIẾT BỊ DẠY HỌC VÀ HỌC LIỆU}

\begin{itemize}[leftmargin=1.2cm]
    \item \textbf{Giáo viên:}
    \begin{itemize}[leftmargin=0.6cm]
        \item Kế hoạch bài dạy chi tiết, bài trình chiếu (PowerPoint/Canva) tích hợp hình ảnh động.
        \item File mô phỏng GeoGebra trực quan hóa đường tròn lượng giác và chuyển động của điểm biểu diễn.
        \item Máy chiếu, máy tính kết nối Internet, phần mềm giả lập máy tính cầm tay Casio fx-580VN X.
        \item Hệ thống Phiếu học tập phân tầng (Worksheet 1, 2, 3) và Rubric đánh giá.
    \end{itemize}
    \item \textbf{Học sinh:}
    \begin{itemize}[leftmargin=0.6cm]
        \item Sách giáo khoa Toán 11, vở ghi bài, giấy nháp.
        \item Thước kẻ, compa, bút màu để vẽ hệ trục và đường tròn lượng giác.
        \item Máy tính cầm tay Casio fx-580VN X (đã kiểm tra pin và chuyển chế độ hiển thị chuẩn).
    \end{itemize}
\end{itemize}

\section*{III. TIẾN TRÌNH DẠY HỌC}

% =========================================================================
% TIẾT 1
% =========================================================================
\subsection*{TIẾT 1. GÓC LƯỢNG GIÁC, ĐƠN VỊ RADIAN VÀ ĐỘ DÀI CUNG TRÒN}
\textit{(Tiết 1 theo PPCT | Tuần 1)}

\subsubsection*{1. Hoạt động 1: Khởi động (Mở đầu)}
\textbf{a) Mục tiêu:} Kích thích sự tò mò của học sinh về sự cần thiết phải mở rộng khái niệm góc hình học sang góc có chiều quay và số đo vượt quá $360^\circ$ hoặc nhận giá trị âm.

\textbf{b) Nội dung:} Giáo viên chiếu video ngắn hoặc hình ảnh chuyển động của kim phút và kim giờ trên mặt đồng hồ; chuyển động của bánh xe đạp khi tiến và lùi. Học sinh trả lời câu hỏi:
\begin{itemize}[leftmargin=0.6cm]
    \item Kim phút quay từ vị trí chỉ số 12 đến số 3 thì quét được góc bao nhiêu độ?
    \item Nếu kim phút quay trọn vẹn 2 vòng rồi đến số 6 thì góc quay tính bằng bao nhiêu độ? Góc hình học lớp 6--10 có biểu diễn được điều này không?
\end{itemize}

\textbf{c) Sản phẩm:} Câu trả lời của học sinh:
\begin{itemize}[leftmargin=0.6cm]
    \item Từ số 12 đến số 3: kim phút quay góc $90^\circ$ theo chiều kim đồng hồ.
    \item Quay 2 vòng rồi đến số 6: kim phút quay $2 \cdot 360^\circ + 180^\circ = 900^\circ$. Góc hình học chỉ từ $0^\circ$ đến $180^\circ$, không thể hiện được chiều quay và số đo lớn hơn $180^\circ$ hay lớn hơn $360^\circ$.
\end{itemize}

\textbf{d) Tổ chức thực hiện:}
\begin{itemize}[leftmargin=0.6cm]
    \item \textit{Chuyển giao nhiệm vụ:} Giáo viên chia lớp thành 35 cặp đôi, trình chiếu hình ảnh đồng hồ cơ và đặt câu hỏi mở đầu.
    \item \textit{Thực hiện nhiệm vụ (Có giàn giáo cho HS TB-Yếu):}
    \begin{scaffoldbox}{Giàn giáo sư phạm 1.1: Hướng dẫn quan sát chiều quay}
        \textbf{Câu hỏi định hướng:}
        \begin{enumerate}[leftmargin=0.5cm]
            \item Mặt đồng hồ chia làm 12 khoảng bằng nhau. Mỗi khoảng tương ứng bao nhiêu độ? ($\dfrac{360^\circ}{12} = 30^\circ$).
            \item Kim đồng hồ quay theo chiều nào? Trong thực tế khi vặn ốc, tiến bánh xe, ta thấy có 2 chiều ngược nhau. Vậy trong Toán học cần quy ước chiều âm, chiều dương như thế nào?
        \end{enumerate}
    \end{scaffoldbox}
    Học sinh quan sát, trao đổi nhanh trong 2 phút.
    \item \textit{Báo cáo, thảo luận:} Đại diện 2 cặp đôi trả lời, các nhóm khác lắng nghe và nhận xét.
    \item \textit{Kết luận, nhận định:} Giáo viên chốt: Cần mở rộng khái niệm góc thành \textbf{Góc lượng giác} có gắn với chiều quay và số đo bất kỳ.
\end{itemize}

\subsubsection*{2. Hoạt động 2: Hình thành kiến thức}

\paragraph{Nội dung 1: Khái niệm góc lượng giác và số đo của góc lượng giác}
\textbf{a) Mục tiêu:} Nhận biết được khái niệm góc lượng giác, tia đầu, tia cuối, chiều quay và quy ước dấu; hiểu công thức tổng quát số đo của góc lượng giác; nhận biết hệ thức Chasles.

\textbf{b) Nội dung:}
\begin{itemize}[leftmargin=0.6cm]
    \item Quy ước chiều quay: Chiều ngược chiều quay của kim đồng hồ là \textbf{chiều dương}; chiều cùng chiều quay của kim đồng hồ là \textbf{chiều âm}.
    \item Khái niệm: Cho hai tia $Ou, Ov$. Khi tia $Om$ quay quanh gốc $O$ từ vị trí tia $Ou$ đến vị trí tia $Ov$, ta nói tia $Om$ quét nên một \textbf{góc lượng giác} có tia đầu là $Ou$, tia cuối là $Ov$, ký hiệu là $(Ou, Ov)$.
    \item Số đo của góc lượng giác: Khi tia $Om$ quay góc $\alpha^\circ$ từ $Ou$ đến $Ov$, số đo các góc lượng giác $(Ou, Ov)$ có dạng:
    $$\text{sđ}(Ou, Ov) = \alpha^\circ + k \cdot 360^\circ \quad (k \in \mathbb{Z})$$
    \item Hệ thức Chasles: Với ba tia $Ou, Ov, Ow$ bất kỳ, ta có:
    $$\text{sđ}(Ou, Ov) + \text{sđ}(Ov, Ow) = \text{sđ}(Ou, Ow) + k \cdot 360^\circ \quad (k \in \mathbb{Z})$$
\end{itemize}

\textbf{c) Sản phẩm:} Học sinh ghi chép chính xác định nghĩa, phân biệt được góc dương, góc âm và xác định được số đo các góc lượng giác cụ thể. Ví dụ: Nếu tia $Om$ quay theo chiều dương $1$ vòng và thêm $45^\circ$ thì $\text{sđ}(Ou, Ov) = 45^\circ + 1 \cdot 360^\circ = 405^\circ$. Nếu quay ngược chiều dương (chiều âm) nửa vòng thì số đo là $-180^\circ$.

\textbf{d) Tổ chức thực hiện:}
\begin{itemize}[leftmargin=0.6cm]
    \item \textit{Chuyển giao nhiệm vụ:} Giáo viên vẽ hình mô phỏng chuyển động quay của tia $Om$ trên bảng hoặc slide trình chiếu, yêu cầu học sinh làm việc cá nhân và trả lời câu hỏi tìm số đo.
    \item \textit{Thực hiện nhiệm vụ (Có giàn giáo cho HS TB-Yếu):}
    \begin{scaffoldbox}{Giàn giáo sư phạm 1.2: Nhớ quy ước dấu}
        \textbf{Mẹo ghi nhớ:}
        \begin{itemize}[leftmargin=0.5cm]
            \item Ngược chiều kim đồng hồ $\implies$ \textbf{Dấu CỘNG} ($+$).
            \item Cùng chiều kim đồng hồ $\implies$ \textbf{Dấu TRỪ} ($-$).
            \item Cứ quay đủ 1 vòng $\implies$ cộng thêm hoặc trừ đi $360^\circ$ (tương ứng bội $k \cdot 360^\circ$).
        \end{itemize}
    \end{scaffoldbox}
    \item \textit{Báo cáo, thảo luận:} Gọi 2 học sinh trung bình lên bảng xác định số đo góc lượng giác khi quay cùng chiều kim đồng hồ $120^\circ$ (kết quả: $-120^\circ + k \cdot 360^\circ$) và quay ngược chiều kim đồng hồ 2 vòng rưỡi (kết quả: $2{,}5 \cdot 360^\circ = 900^\circ$).
    \item \textit{Kết luận, nhận định:} Giáo viên chuẩn hóa kiến thức, nhấn mạnh: Hai tia $Ou, Ov$ cho trước tạo ra \textit{vô số} góc lượng giác sai khác nhau một bội số nguyên của $360^\circ$.
\end{itemize}

\paragraph{Nội dung 2: Đơn vị đo góc Radian và độ dài cung tròn}
\textbf{a) Mục tiêu:} Nhận biết định nghĩa radian; chuyển đổi thành thạo giữa đơn vị độ và radian; tính được độ dài cung tròn $l = \alpha R$.

\textbf{b) Nội dung:}
\begin{itemize}[leftmargin=0.6cm]
    \item Định nghĩa: Radian là số đo của một góc ở tâm chắn một cung có độ dài bằng bán kính của đường tròn đó. Ký hiệu là rad.
    \item Mối quan hệ giữa độ và radian:
    $$180^\circ = \pi \text{ rad} \implies 1^\circ = \dfrac{\pi}{180} \text{ rad}, \quad 1 \text{ rad} = \left(\dfrac{180}{\pi}\right)^\circ \approx 57^\circ 17' 45''$$
    \item Công thức chuyển đổi:
    $$\alpha \text{ rad} = a^\circ \cdot \dfrac{\pi}{180^\circ}, \quad a^\circ = \alpha \cdot \dfrac{180^\circ}{\pi}$$
    \item Công thức độ dài cung tròn: Trên đường tròn bán kính $R$, cung có số đo $\alpha$ rad có độ dài là:
    $$l = \alpha \cdot R \quad (\alpha \text{ tính bằng rad})$$
\end{itemize}

\textbf{c) Sản phẩm:} Bảng chuyển đổi các góc đặc biệt thường gặp:
\begin{center}
\begin{tabular}{lcccccccc}
\toprule
\textbf{Độ} & $0^\circ$ & $30^\circ$ & $45^\circ$ & $60^\circ$ & $90^\circ$ & $120^\circ$ & $135^\circ$ & $180^\circ$ \\
\midrule
\textbf{Radian} & $0$ & $\dfrac{\pi}{6}$ & $\dfrac{\pi}{4}$ & $\dfrac{\pi}{3}$ & $\dfrac{\pi}{2}$ & $\dfrac{2\pi}{3}$ & $\dfrac{3\pi}{4}$ & $\pi$ \\
\bottomrule
\end{tabular}
\end{center}

\textbf{d) Tổ chức thực hiện:}
\begin{itemize}[leftmargin=0.6cm]
    \item \textit{Chuyển giao nhiệm vụ:} Giáo viên giao bài tập điền vào chỗ trống bảng góc đặc biệt và tính độ dài cung tròn chắn góc $\dfrac{\pi}{4}$ trên đường tròn bán kính $R = 10\text{ cm}$.
    \item \textit{Thực hiện nhiệm vụ (Năng lực số 5.2NC1b):}
    \begin{aibox}{Thực hành Năng lực số: Máy tính cầm tay Casio fx-580VN X}
        \textbf{Quy trình chuyển đổi đơn vị góc trên MTCT:}
        \begin{itemize}[leftmargin=0.5cm]
            \item Cài đặt đơn vị đo Radian: Nhấn \texttt{SHIFT} $\to$ \texttt{MENU} $\to$ \texttt{2} (Angle Unit) $\to$ Chọn \texttt{2} (Radian). Màn hình xuất hiện chữ \textbf{R}.
            \item Đổi từ Độ sang Radian: Nhập số đo độ $\to$ Nhấn \texttt{OPTN} $\to$ \texttt{2} (Angle Unit) $\to$ \texttt{1} ($^\circ$) $\to$ Nhấn \texttt{=}.\\
            \textit{Ví dụ:} Nhập $60^\circ \to$ Máy hiện $\dfrac{1}{3}\pi$.
        \end{itemize}
    \end{aibox}
    \item \textit{Báo cáo, thảo luận:} Học sinh thực hành bấm máy tại chỗ trên 35 MTCT, 1 học sinh đọc kết quả độ dài cung: $l = \dfrac{\pi}{4} \cdot 10 = \dfrac{5\pi}{2} \approx 7{,}85\text{ cm}$.
    \item \textit{Kết luận, nhận định:} Giáo viên đánh giá kỹ năng thao tác máy tính của học sinh, lưu ý khi tính độ dài cung tròn $l = \alpha R$, số đo $\alpha$ \textbf{bắt buộc phải đổi sang radian}.
\end{itemize}

\subsubsection*{3. Hoạt động 3: Luyện tập}
\textbf{a) Mục tiêu:} Củng cố kỹ năng đổi số đo giữa độ và radian; tính số đo góc lượng giác; tính độ dài cung tròn.

\textbf{b) Nội dung:} Học sinh hoàn thành Phiếu học tập số 1 (Worksheet 1) gồm 3 bài tập:
\begin{itemize}[leftmargin=0.6cm]
    \item \textbf{Bài 1:} Đổi các số đo sau sang radian: $75^\circ; -150^\circ; 225^\circ$.
    \item \textbf{Bài 2:} Đổi các số đo sau sang độ: $\dfrac{3\pi}{8}; -\dfrac{5\pi}{6}; 2\text{ rad}$.
    \item \textbf{Bài 3:} Một bánh xe máy có bán kính $R = 25\text{ cm}$. Khi bánh xe quay được góc $\alpha = \dfrac{7\pi}{3}\text{ rad}$, bánh xe đã lăn được một quãng đường dài bao nhiêu cm?
\end{itemize}

\textbf{c) Sản phẩm:} Lời giải chi tiết của học sinh:
\begin{itemize}[leftmargin=0.6cm]
    \item \textbf{Bài 1:} $75^\circ = 75 \cdot \dfrac{\pi}{180} = \dfrac{5\pi}{12}\text{ rad}$; $-150^\circ = -150 \cdot \dfrac{\pi}{180} = -\dfrac{5\pi}{6}\text{ rad}$; $225^\circ = \dfrac{5\pi}{4}\text{ rad}$.
    \item \textbf{Bài 2:} $\dfrac{3\pi}{8}\text{ rad} = \dfrac{3 \cdot 180^\circ}{8} = 67{,}5^\circ = 67^\circ 30'$; $-\dfrac{5\pi}{6}\text{ rad} = -150^\circ$; $2\text{ rad} = \left(\dfrac{2 \cdot 180}{\pi}\right)^\circ \approx 114^\circ 35'$.
    \item \textbf{Bài 3:} Quãng đường bánh xe lăn được bằng độ dài cung: $l = \alpha \cdot R = \dfrac{7\pi}{3} \cdot 25 = \dfrac{175\pi}{3} \approx 183{,}26\text{ cm}$.
\end{itemize}

\textbf{d) Tổ chức thực hiện:}
\begin{itemize}[leftmargin=0.6cm]
    \item \textit{Chuyển giao nhiệm vụ:} Giáo viên phát Phiếu học tập số 1, yêu cầu làm việc độc lập 5 phút, sau đó đổi bài chéo cặp đôi kiểm tra.
    \item \textit{Thực hiện nhiệm vụ:} Học sinh thực hiện tính toán và kiểm tra kết quả bằng MTCT. Giáo viên đi quanh lớp quan sát, hỗ trợ các bạn yếu đổi phân số.
    \item \textit{Báo cáo, thảo luận:} 3 học sinh lên bảng trình bày 3 bài. Lớp nhận xét chéo.
    \item \textit{Kết luận, nhận định:} Giáo viên nhận xét, chốt điểm cần lưu ý khi làm tròn số và kiểm tra đơn vị hiển thị trên máy tính cầm tay.
\end{itemize}

\subsubsection*{4. Hoạt động 4: Vận dụng}
\textbf{a) Mục tiêu:} Vận dụng công thức độ dài cung tròn và tốc độ góc vào bài toán thực tiễn chuyển động kim đồng hồ.

\textbf{b) Nội dung:} Kim phút của một đồng hồ treo tường có chiều dài $r = 12\text{ cm}$. Trong khoảng thời gian 25 phút, đầu mút của kim phút vạch nên một cung tròn có độ dài xấp xỉ bằng bao nhiêu centimet?

\textbf{c) Sản phẩm:} Lời giải:
\begin{itemize}[leftmargin=0.6cm]
    \item Trong 60 phút, kim phút quay trọn 1 vòng tương ứng góc $2\pi\text{ rad}$.
    \item Do đó, trong 25 phút, góc quay của kim phút (xét về độ lớn) là:
    $$\alpha = 25 \cdot \dfrac{2\pi}{60} = \dfrac{5\pi}{6}\text{ rad}$$
    \item Độ dài cung tròn mà đầu mút kim phút vạch nên là:
    $$l = \alpha \cdot r = \dfrac{5\pi}{6} \cdot 12 = 10\pi \approx 31{,}42\text{ cm}$$
\end{itemize}

\textbf{d) Tổ chức thực hiện:}
\begin{itemize}[leftmargin=0.6cm]
    \item \textit{Chuyển giao nhiệm vụ:} Giáo viên giao bài toán vận dụng, yêu cầu học sinh thảo luận cặp đôi trong 3 phút.
    \item \textit{Thực hiện nhiệm vụ:} Học sinh xác định tỉ số thời gian $\dfrac{25}{60}$ để tính góc $\alpha$, sau đó áp dụng $l = \alpha r$.
    \item \textit{Báo cáo, thảo luận:} 1 nhóm xung phong trình bày cách tính. Nhóm khác nêu cách quy đổi ra độ trước ($25\text{ phút} = 150^\circ$) rồi đổi ra radian.
    \item \textit{Kết luận, nhận định:} Giáo viên tổng kết Tiết 1, dặn dò học sinh ghi nhớ bảng chuyển đổi góc đặc biệt và xem trước khái niệm Đường tròn lượng giác.
\end{itemize}

% =========================================================================
% TIẾT 2
% =========================================================================
\vspace{10pt}
\subsection*{TIẾT 2. ĐƯỜNG TRÒN LƯỢNG GIÁC VÀ GIÁ TRỊ LƯỢNG GIÁC CỦA GÓC LƯỢNG GIÁC}
\textit{(Tiết 2 theo PPCT | Tuần 1)}

\subsubsection*{1. Hoạt động 1: Khởi động (Mở đầu)}
\textbf{a) Mục tiêu:} Tái hiện mối liên hệ giữa tọa độ điểm trên hệ trục tọa độ và giá trị $\sin, \cos$ đã học ở lớp 10 (trong nửa đường tròn đơn vị), mở rộng ra toàn bộ đường tròn.

\textbf{b) Nội dung:} Giáo viên chiếu hình ảnh vòng quay Mặt Trời (Ferris Wheel). Một cabin gắn tại điểm $M$ quay quanh tâm $O$. Khi gắn hệ trục tọa độ $Oxy$ với gốc tọa độ đặt tại tâm trục quay:
\begin{itemize}[leftmargin=0.6cm]
    \item Vị trí cabin $M$ được xác định bởi tọa độ $(x; y)$ như thế nào?
    \item Ở lớp 10, với góc $\alpha \in [0^\circ; 180^\circ]$, $\cos\alpha$ và $\sin\alpha$ là gì của điểm $M$? Khi cabin quay xuống phía dưới trục hoành (góc $> 180^\circ$), tung độ $y$ âm thì $\sin\alpha$ nhận giá trị thế nào?
\end{itemize}

\textbf{c) Sản phẩm:} Học sinh nhận ra: Ở lớp 10, trên nửa đường tròn đơn vị, $\cos\alpha = x_M$ và $\sin\alpha = y_M$. Khi mở rộng ra cả đường tròn, với góc tùy ý, ta hoàn toàn có thể định nghĩa tương tự: hoành độ là $\cos$, tung độ là $\sin$.

\textbf{d) Tổ chức thực hiện:}
\begin{itemize}[leftmargin=0.6cm]
    \item \textit{Chuyển giao nhiệm vụ:} Giáo viên đặt câu hỏi gợi mở kết hợp hình ảnh động GeoGebra vòng quay.
    \item \textit{Thực hiện nhiệm vụ:} Học sinh suy nghĩ, thảo luận nhanh với bạn cùng bàn.
    \item \textit{Báo cáo, thảo luận:} 2 học sinh phát biểu liên hệ kiến thức lớp 10.
    \item \textit{Kết luận, nhận định:} Giáo viên dẫn dắt vào bài mới: Thiết lập khái niệm Đường tròn lượng giác và định nghĩa 4 giá trị lượng giác của góc bất kỳ.
\end{itemize}

\subsubsection*{2. Hoạt động 2: Hình thành kiến thức}

\paragraph{Nội dung 1: Đường tròn lượng giác và Điểm biểu diễn của góc lượng giác}
\textbf{a) Mục tiêu:} Nhận biết khái niệm đường tròn lượng giác; biết cách xác định điểm biểu diễn của góc lượng giác có số đo $\alpha$.

\textbf{b) Nội dung:}
\begin{itemize}[leftmargin=0.6cm]
    \item \textbf{Định nghĩa:} Trong mặt phẳng tọa độ $Oxy$, đường tròn lượng giác là đường tròn tâm $O$, bán kính $R = 1$, trên đó đã chọn chiều ngược chiều kim đồng hồ làm chiều dương và lấy điểm $A(1; 0)$ làm gốc của đường tròn.
    \item \textbf{Điểm biểu diễn:} Điểm $M$ trên đường tròn lượng giác được gọi là điểm biểu diễn của góc lượng giác có số đo $\alpha$ nếu số đo của góc lượng giác $(OA, OM)$ bằng $\alpha$, tức là $\text{sđ}(OA, OM) = \alpha$.
\end{itemize}

\textbf{c) Sản phẩm:} Học sinh vẽ đúng đường tròn lượng giác và biểu diễn được các điểm ứng với góc đặc biệt:
\begin{itemize}[leftmargin=0.6cm]
    \item $\alpha = 0 \implies M \equiv A(1; 0)$.
    \item $\alpha = \dfrac{\pi}{2} = 90^\circ \implies M \equiv B(0; 1)$.
    \item $\alpha = \pi = 180^\circ \implies M \equiv A'(-1; 0)$.
    \item $\alpha = \dfrac{3\pi}{2} = 270^\circ = -90^\circ \implies M \equiv B'(0; -1)$.
    \item $\alpha = \dfrac{\pi}{4} = 45^\circ \implies M$ là điểm chính giữa cung $AB$.
\end{itemize}

\textbf{d) Tổ chức thực hiện:}
\begin{itemize}[leftmargin=0.6cm]
    \item \textit{Chuyển giao nhiệm vụ:} Giáo viên hướng dẫn học sinh dùng compa vẽ hệ trục tọa độ $Oxy$ và đường tròn bán kính $R = 1$ (lấy 2 ô li tập làm 1 đơn vị). Yêu cầu xác định các điểm $A, B, A', B'$ và biểu diễn điểm $M$ với góc $\alpha = \dfrac{5\pi}{6}$.
    \item \textit{Thực hiện nhiệm vụ (Giàn giáo sư phạm):}
    \begin{scaffoldbox}{Giàn giáo sư phạm 2.1: Cách xác định điểm biểu diễn góc $\alpha$}
        \textbf{Quy tắc 3 bước:}
        \begin{enumerate}[leftmargin=0.5cm]
            \item Xuất phát luôn luôn từ điểm gốc $A(1; 0)$.
            \item Xác định chiều quay: Nếu $\alpha > 0$ quay ngược chiều KĐH; nếu $\alpha < 0$ quay cùng chiều KĐH.
            \item Chia góc để xác định vị trí: Ví dụ $\dfrac{5\pi}{6} = \pi - \dfrac{\pi}{6}$, quay đến nửa đường tròn (điểm $A'$) rồi lùi lại $\dfrac{\pi}{6} = 30^\circ$ theo chiều âm.
        \end{enumerate}
    \end{scaffoldbox}
    \item \textit{Báo cáo, thảo luận:} Giáo viên gọi 1 học sinh lên bảng biểu diễn điểm ứng với góc $-\dfrac{\pi}{3}$. Lớp quan sát, nhận xét.
    \item \textit{Kết luận, nhận định:} Giáo viên nhấn mạnh: Mỗi góc lượng giác $\alpha$ xác định duy nhất một điểm $M$ trên đường tròn lượng giác. Ngược lại, một điểm $M$ biểu diễn vô số góc có số đo sai khác $k2\pi$.
\end{itemize}

\paragraph{Nội dung 2: Khái niệm giá trị lượng giác và Dấu của các giá trị lượng giác}
\textbf{a) Mục tiêu:} Nhận biết và viết đúng định nghĩa $\sin\alpha, \cos\alpha, \tan\alpha, \cot\alpha$; xác định chính xác dấu của 4 giá trị lượng giác trong 4 góc phần tư; mô tả bảng giá trị góc đặc biệt.

\textbf{b) Nội dung:}
\begin{itemize}[leftmargin=0.6cm]
    \item \textbf{Định nghĩa:} Giả sử điểm $M(x; y)$ là điểm biểu diễn của góc lượng giác $\alpha$ trên đường tròn lượng giác. Khi đó:
    \begin{itemize}[leftmargin=0.5cm]
        \item $\cos\alpha = x$ (Hoành độ của điểm $M$ -- trục $Ox$ là trục cosin).
        \item $\sin\alpha = y$ (Tung độ của điểm $M$ -- trục $Oy$ là trục sin).
        \item $\tan\alpha = \dfrac{\sin\alpha}{\cos\alpha} = \dfrac{y}{x} \quad (\alpha \neq \dfrac{\pi}{2} + k\pi, k \in \mathbb{Z})$.
        \item $\cot\alpha = \dfrac{\cos\alpha}{\sin\alpha} = \dfrac{x}{y} \quad (\alpha \neq k\pi, k \in \mathbb{Z})$.
    \end{itemize}
    \item \textbf{Ý nghĩa giới hạn:} Vì điểm $M$ thuộc đường tròn đơn vị nên:
    $$-1 \le \cos\alpha \le 1, \quad -1 \le \sin\alpha \le 1 \quad (\forall \alpha \in \mathbb{R})$$
    \item \textbf{Bảng xét dấu trên 4 góc phần tư:}
    \begin{center}
    \begin{tabular}{ccccc}
    \toprule
    \textbf{Giá trị lượng giác} & \textbf{Góc phần tư I} & \textbf{Góc phần tư II} & \textbf{Góc phần tư III} & \textbf{Góc phần tư IV} \\
    & $\left(0 < \alpha < \dfrac{\pi}{2}\right)$ & $\left(\dfrac{\pi}{2} < \alpha < \pi\right)$ & $\left(\pi < \alpha < \dfrac{3\pi}{2}\right)$ & $\left(\dfrac{3\pi}{2} < \alpha < 2\pi\right)$ \\
    \midrule
    $\sin\alpha$ & $+$ & $+$ & $-$ & $-$ \\
    $\cos\alpha$ & $+$ & $-$ & $-$ & $+$ \\
    $\tan\alpha$ & $+$ & $-$ & $+$ & $-$ \\
    $\cot\alpha$ & $+$ & $-$ & $+$ & $-$ \\
    \bottomrule
    \end{tabular}
    \end{center}
\end{itemize}

\textbf{c) Sản phẩm:} Học sinh lập được bảng xét dấu, giải thích được nguyên nhân dấu của từng góc phần tư dựa vào dấu của hoành độ $x$ và tung độ $y$ của điểm $M$.

\textbf{d) Tổ chức thực hiện:}
\begin{itemize}[leftmargin=0.6cm]
    \item \textit{Chuyển giao nhiệm vụ:} Giáo viên chia lớp theo nhóm 4 học sinh, yêu cầu dựa vào vị trí của điểm $M(x; y)$ trong từng góc phần tư để điền dấu $+$ hoặc $-$ vào bảng dấu.
    \item \textit{Thực hiện nhiệm vụ (Giàn giáo cho HS TB-Yếu):}
    \begin{scaffoldbox}{Giàn giáo sư phạm 2.2: Câu khẩu quyết nhớ dấu lượng giác}
        \textbf{Khẩu quyết dân gian:} \textit{''Nhất cả, nhị sin, tam tang, tứ cos''}
        \begin{itemize}[leftmargin=0.5cm]
            \item \textbf{Nhất cả:} Góc phần tư I $\implies$ Tất cả đều DƯƠNG ($+$).
            \item \textbf{Nhị sin:} Góc phần tư II $\implies$ Chỉ có $\sin$ DƯƠNG ($+$), còn lại âm.
            \item \textbf{Tam tang:} Góc phần tư III $\implies$ Chỉ có $\tan$ và $\cot$ DƯƠNG ($+$), còn lại âm.
            \item \textbf{Tứ cos:} Góc phần tư IV $\implies$ Chỉ có $\cos$ DƯƠNG ($+$), còn lại âm.
        \end{itemize}
    \end{scaffoldbox}
    \item \textit{Tích hợp Năng lực AI (Mã 11.C3.1):}
    \begin{aibox}{Thực hành Năng lực AI: Kỹ thuật Prompt kiểm chứng bảng dấu}
        \textbf{Mẫu câu lệnh (Prompt):}\\
        \texttt{"Hãy giải thích dấu của $\sin\alpha, \cos\alpha, \tan\alpha$ khi góc $\alpha$ nằm ở góc phần tư thứ III bằng cách phân tích tọa độ hoành độ $x$ và tung độ $y$ trên đường tròn lượng giác đơn vị."}\\
        $\implies$ Học sinh nhận diện: AI phản hồi giải thích chính xác vì ở góc phần tư III, cả $x < 0$ và $y < 0$ nên $\cos < 0, \sin < 0$, dẫn tới $\tan = \dfrac{y}{x} > 0$.
    \end{aibox}
    \item \textit{Báo cáo, thảo luận:} Đại diện 1 nhóm báo cáo bảng dấu hoàn chỉnh, lớp nhận xét.
    \item \textit{Kết luận, nhận định:} Giáo viên chốt bảng dấu và bảng giá trị góc đặc biệt (học sinh tra cứu SGK và luyện bấm máy tính cầm tay).
\end{itemize}

\subsubsection*{3. Hoạt động 3: Luyện tập}
\textbf{a) Mục tiêu:} Rèn luyện kỹ năng xác định dấu; tính giá trị lượng giác bằng MTCT Casio fx-580VN X; tìm góc lượng giác khi biết giá trị.

\textbf{b) Nội dung:} Học sinh thực hiện các bài tập trên Phiếu học tập số 2 (Worksheet 2):
\begin{itemize}[leftmargin=0.6cm]
    \item \textbf{Bài 1:} Xác định dấu của: $\sin 156^\circ; \cos\left(-\dfrac{\pi}{3}\right); \tan\dfrac{4\pi}{3}; \cot\left(-215^\circ\right)$.
    \item \textbf{Bài 2:} Sử dụng MTCT tính giá trị đúng hoặc gần đúng (làm tròn 4 chữ số thập phân):
    $$\sin\left(-\dfrac{2\pi}{7}\right); \quad \cos 215^\circ; \quad \tan\dfrac{5\pi}{12}; \quad \cot\left(\dfrac{3\pi}{5}\right)$$
    \item \textbf{Bài 3:} Tìm số đo góc $\alpha \in [0; \pi]$ biết $\cos\alpha = -\dfrac{1}{2}$.
\end{itemize}

\textbf{c) Sản phẩm:} Lời giải:
\begin{itemize}[leftmargin=0.6cm]
    \item \textbf{Bài 1:}
    \begin{itemize}[leftmargin=0.5cm]
        \item $156^\circ \in (90^\circ; 180^\circ)$ (Góc phần tư II) $\implies \sin 156^\circ > 0$.
        \item $-\dfrac{\pi}{3} = -60^\circ$ (Góc phần tư IV) $\implies \cos\left(-\dfrac{\pi}{3}\right) > 0$.
        \item $\dfrac{4\pi}{3} = \pi + \dfrac{\pi}{3}$ (Góc phần tư III) $\implies \tan\dfrac{4\pi}{3} > 0$.
        \item $-215^\circ = -180^\circ - 35^\circ$ (Góc phần tư II) $\implies \cot\left(-215^\circ\right) < 0$.
    \end{itemize}
    \item \textbf{Bài 2:} Bấm máy MTCT Casio fx-580VN X:
    \begin{itemize}[leftmargin=0.5cm]
        \item $\sin\left(-\dfrac{2\pi}{7}\right) \approx -0{,}7818$.
        \item $\cos 215^\circ \approx -0{,}8192$.
        \item $\tan\dfrac{5\pi}{12} = 2 + \sqrt{3} \approx 3{,}7321$.
        \item $\cot\left(\dfrac{3\pi}{5}\right) = \dfrac{1}{\tan\left(\dfrac{3\pi}{5}\right)} \approx -0{,}3249$.
    \end{itemize}
    \item \textbf{Bài 3:} Trên MTCT: Nhấn \texttt{SHIFT} $\to$ \texttt{cos} $\to$ Nhập $-\dfrac{1}{2} \to$ Kết quả: $\alpha = \dfrac{2\pi}{3} = 120^\circ$.
\end{itemize}

\textbf{d) Tổ chức thực hiện:}
\begin{itemize}[leftmargin=0.6cm]
    \item \textit{Chuyển giao nhiệm vụ:} Học sinh làm việc cá nhân trong 7 phút. Giáo viên hướng dẫn riêng cách nhập hàm $\cot$ trên MTCT (nhấn $1 \div \tan$).
    \item \textit{Thực hiện nhiệm vụ:} Học sinh làm bài, đối chiếu kết quả theo cặp.
    \item \textit{Báo cáo, thảo luận:} 2 học sinh lên bảng ghi kết quả, chỉ rõ thao tác phím bấm máy tính.
    \item \textit{Kết luận, nhận định:} Giáo viên chuẩn hóa thao tác máy tính, nhắc nhở lỗi sai kinh điển: quên chuyển đổi giữa chế độ độ (D) và radian (R) khi bấm máy.
\end{itemize}

\subsubsection*{4. Hoạt động 4: Vận dụng}
\textbf{a) Mục tiêu:} Ứng dụng giá trị lượng giác để tính tọa độ và độ cao trong chuyển động thực tế của vòng quay đu quay.

\textbf{b) Nội dung:} Một vòng đu quay Mặt Trời có bán kính $R = 40\text{ m}$, tâm đu quay ở độ cao $50\text{ m}$ so với mặt đất. Giả sử vòng quay quay đều theo chiều ngược chiều kim đồng hồ với vận tốc 1 vòng mất 12 phút. Ban đầu cabin ở vị trí thấp nhất (điểm $B'$ có độ cao $10\text{ m}$). Hỏi sau 4 phút, cabin ở độ cao bao nhiêu mét so với mặt đất?

\textbf{c) Sản phẩm:} Lời giải:
\begin{itemize}[leftmargin=0.6cm]
    \item Vận tốc góc của đu quay: $\omega = \dfrac{2\pi}{12} = \dfrac{\pi}{6}\text{ (rad/phút)}$.
    \item Sau $t = 4\text{ phút}$, góc quay được là: $\alpha = 4 \cdot \dfrac{\pi}{6} = \dfrac{2\pi}{3}\text{ rad} = 120^\circ$.
    \item Vị trí ban đầu ở điểm $B'$ tương ứng góc lượng giác $-90^\circ$ (hoặc $\dfrac{3\pi}{2}$). Sau khi quay $120^\circ$ theo chiều dương, vị trí của cabin tương ứng với góc lượng giác:
    $$\beta = -90^\circ + 120^\circ = 30^\circ$$
    \item Tung độ của cabin so với tâm đu quay là: $y = R \cdot \sin 30^\circ = 40 \cdot \dfrac{1}{2} = 20\text{ m}$.
    \item Độ cao của cabin so với mặt đất là: $h = 50 + y = 50 + 20 = 70\text{ m}$.
\end{itemize}

\textbf{d) Tổ chức thực hiện:}
\begin{itemize}[leftmargin=0.6cm]
    \item \textit{Chuyển giao nhiệm vụ:} Giáo viên giao bài toán cho các nhóm 4 người, hướng dẫn phân tích hình học.
    \item \textit{Thực hiện nhiệm vụ:} Nhóm thảo luận, vẽ sơ đồ hình tròn tọa độ tâm $O(0; 50)$ bán kính $R = 40$.
    \item \textit{Báo cáo, thảo luận:} 1 nhóm đại diện lên bảng vẽ hình và trình bày lời giải.
    \item \textit{Kết luận, nhận định:} Giáo viên nhận xét bài làm, nhấn mạnh tính ứng dụng sâu sắc của đường tròn lượng giác trong chuyển động chu kỳ của Vật lý và Kỹ thuật.
\end{itemize}

% =========================================================================
% TIẾT 3
% =========================================================================
\vspace{10pt}
\subsection*{TIẾT 3. CÁC HỆ THỨC LƯỢNG GIÁC CƠ BẢN VÀ GIÁ TRỊ LƯỢNG GIÁC CỦA CÁC GÓC LIÊN QUAN ĐẶC BIỆT}
\textit{(Tiết 3 theo PPCT | Tuần 1)}

\subsubsection*{1. Hoạt động 1: Khởi động (Mở đầu)}
\textbf{a) Mục tiêu:} Gợi nhớ định lý Pythagore và hệ thức $\sin^2\alpha + \cos^2\alpha = 1$ ở tam giác vuông lớp 9 và nửa đường tròn đơn vị lớp 10; tạo tình huống xuất phát mở rộng cho góc lượng giác bất kỳ.

\textbf{b) Nội dung:} Cho góc lượng giác $\alpha$ bất kỳ với điểm biểu diễn $M(x; y)$ trên đường tròn lượng giác.
\begin{itemize}[leftmargin=0.6cm]
    \item Hãy viết phương trình đường tròn lượng giác tâm $O(0; 0)$ bán kính $R = 1$.
    \item Từ đó suy ra hệ thức liên hệ giữa hoành độ $x = \cos\alpha$ và tung độ $y = \sin\alpha$. Hệ thức này có đúng với mọi góc $\alpha$ không?
\end{itemize}

\textbf{c) Sản phẩm:} Phương trình đường tròn tâm $O$ bán kính $1$ là $x^2 + y^2 = 1$. Do đó thay $x = \cos\alpha, y = \sin\alpha$ ta được:
$$\sin^2\alpha + \cos^2\alpha = 1 \quad (\forall \alpha \in \mathbb{R})$$

\textbf{d) Tổ chức thực hiện:}
\begin{itemize}[leftmargin=0.6cm]
    \item \textit{Chuyển giao nhiệm vụ:} Giáo viên đặt câu hỏi, gọi 1 học sinh trả lời miệng.
    \item \textit{Thực hiện nhiệm vụ:} Học sinh đối chiếu phương trình đường tròn và định nghĩa $\cos, \sin$.
    \item \textit{Báo cáo, thảo luận:} Học sinh nêu công thức và khẳng định luôn đúng với mọi $\alpha$.
    \item \textit{Kết luận, nhận định:} Giáo viên dẫn dắt vào bài: Nghiên cứu trọn vẹn 4 hệ thức lượng giác cơ bản và mối liên hệ giữa các góc có tính chất đối, bù, phụ, hơn kém nhau $\pi$.
\end{itemize}

\subsubsection*{2. Hoạt động 2: Hình thành kiến thức}

\paragraph{Nội dung 1: Các hệ thức lượng giác cơ bản}
\textbf{a) Mục tiêu:} Nhận biết và chứng minh được 4 hệ thức lượng giác cơ bản; biết điều kiện xác định của từng hệ thức.

\textbf{b) Nội dung:}
Với mọi góc lượng giác $\alpha$ thỏa mãn điều kiện xác định, ta có các hệ thức:
\begin{enumerate}[leftmargin=0.8cm]
    \item $\sin^2\alpha + \cos^2\alpha = 1$
    \item $1 + \tan^2\alpha = \dfrac{1}{\cos^2\alpha} \quad \left(\alpha \neq \dfrac{\pi}{2} + k\pi, k \in \mathbb{Z}\right)$
    \item $1 + \cot^2\alpha = \dfrac{1}{\sin^2\alpha} \quad (\alpha \neq k\pi, k \in \mathbb{Z})$
    \item $\tan\alpha \cdot \cot\alpha = 1 \quad \left(\alpha \neq k\dfrac{\pi}{2}, k \in \mathbb{Z}\right)$
\end{enumerate}

\textbf{c) Sản phẩm:} Học sinh trình bày chứng minh:
\begin{itemize}[leftmargin=0.6cm]
    \item Hệ thức 2: $1 + \tan^2\alpha = 1 + \dfrac{\sin^2\alpha}{\cos^2\alpha} = \dfrac{\cos^2\alpha + \sin^2\alpha}{\cos^2\alpha} = \dfrac{1}{\cos^2\alpha}$.
    \item Hệ thức 3: $1 + \cot^2\alpha = 1 + \dfrac{\cos^2\alpha}{\sin^2\alpha} = \dfrac{\sin^2\alpha + \cos^2\alpha}{\sin^2\alpha} = \dfrac{1}{\sin^2\alpha}$.
    \item Hệ thức 4: $\tan\alpha \cdot \cot\alpha = \dfrac{\sin\alpha}{\cos\alpha} \cdot \dfrac{\cos\alpha}{\sin\alpha} = 1$.
\end{itemize}

\textbf{d) Tổ chức thực hiện:}
\begin{itemize}[leftmargin=0.6cm]
    \item \textit{Chuyển giao nhiệm vụ:} Giáo viên yêu cầu học sinh làm việc cá nhân chứng minh hệ thức (2), (3), (4) dựa vào định nghĩa của $\tan\alpha, \cot\alpha$ và hệ thức (1).
    \item \textit{Thực hiện nhiệm vụ (Giàn giáo cho HS TB-Yếu):}
    \begin{scaffoldbox}{Giàn giáo sư phạm 3.1: Quy đồng mẫu số phân thức lượng giác}
        Nhắc lại: $\tan\alpha = \dfrac{\sin\alpha}{\cos\alpha}$. Khi tính $1 + \tan^2\alpha$, ta thay số $1 = \dfrac{\cos^2\alpha}{\cos^2\alpha}$ rồi cộng hai phân số cùng mẫu.
    \end{scaffoldbox}
    \item \textit{Báo cáo, thảo luận:} 2 học sinh lên bảng trình bày phép biến đổi chứng minh. Lớp nhận xét.
    \item \textit{Kết luận, nhận định:} Giáo viên nhấn mạnh: 4 hệ thức này là công cụ cốt lõi để tính toán các giá trị lượng giác khi biết trước 1 giá trị.
\end{itemize}

\paragraph{Nội dung 2: Giá trị lượng giác của các góc có liên quan đặc biệt}
\textbf{a) Mục tiêu:} Nhận biết và mô tả được mối quan hệ lượng giác giữa các góc đối nhau, bù nhau, hơn kém nhau $\pi$, phụ nhau; vận dụng linh hoạt để rút gọn biểu thức.

\textbf{b) Nội dung:}
\begin{itemize}[leftmargin=0.6cm]
    \item \textbf{Hai góc đối nhau ($\alpha$ và $-\alpha$):}
    $$\cos(-\alpha) = \cos\alpha, \quad \sin(-\alpha) = -\sin\alpha, \quad \tan(-\alpha) = -\tan\alpha, \quad \cot(-\alpha) = -\cot\alpha$$
    \item \textbf{Hai góc bù nhau ($\alpha$ và $\pi - \alpha$):}
    $$\sin(\pi - \alpha) = \sin\alpha, \quad \cos(\pi - \alpha) = -\cos\alpha, \quad \tan(\pi - \alpha) = -\tan\alpha, \quad \cot(\pi - \alpha) = -\cot\alpha$$
    \item \textbf{Hai góc hơn kém nhau $\pi$ ($\alpha$ và $\alpha + \pi$):}
    $$\sin(\alpha + \pi) = -\sin\alpha, \quad \cos(\alpha + \pi) = -\cos\alpha, \quad \tan(\alpha + \pi) = \tan\alpha, \quad \cot(\alpha + \pi) = \cot\alpha$$
    \item \textbf{Hai góc phụ nhau $\left(\alpha \text{ và } \dfrac{\pi}{2} - \alpha\right)$:}
    $$\sin\left(\dfrac{\pi}{2} - \alpha\right) = \cos\alpha, \quad \cos\left(\dfrac{\pi}{2} - \alpha\right) = \sin\alpha, \quad \tan\left(\dfrac{\pi}{2} - \alpha\right) = \cot\alpha, \quad \cot\left(\dfrac{\pi}{2} - \alpha\right) = \tan\alpha$$
\end{itemize}

\textbf{c) Sản phẩm:} Học sinh giải thích được tính đối xứng của các điểm biểu diễn trên đường tròn lượng giác:
\begin{itemize}[leftmargin=0.6cm]
    \item Góc $\alpha$ và $-\alpha$: Hai điểm biểu diễn đối xứng nhau qua trục hoành $Ox \implies$ hoành độ bằng nhau ($\cos$ bằng nhau), tung độ đối nhau ($\sin$ đối nhau).
    \item Góc $\alpha$ và $\pi - \alpha$: Hai điểm đối xứng qua trục tung $Oy \implies$ tung độ bằng nhau ($\sin$ bằng nhau), hoành độ đối nhau ($\cos$ đối nhau).
\end{itemize}

\textbf{d) Tổ chức thực hiện:}
\begin{itemize}[leftmargin=0.6cm]
    \item \textit{Chuyển giao nhiệm vụ:} Giáo viên trình chiếu đường tròn lượng giác với các cặp điểm đối xứng qua $Ox, Oy$, yêu cầu học sinh rút ra nhận xét về mối quan hệ giữa các giá trị lượng giác.
    \item \textit{Thực hiện nhiệm vụ (Giàn giáo sư phạm):}
    \begin{scaffoldbox}{Giàn giáo sư phạm 3.2: Thơ vui ghi nhớ công thức cung liên kết}
        \textbf{Bài thơ thần chú:}\\
        \textit{''Cos đối, Sin bù, Phụ chéo, Hơn kém pi tang''}\\
        \textbf{Giải nghĩa:}\\
        - \textbf{Cos đối:} Chỉ có $\cos$ giữ nguyên dấu: $\cos(-\alpha) = \cos\alpha$, các hàm khác đổi dấu.\\
        - \textbf{Sin bù:} Chỉ có $\sin$ giữ nguyên dấu: $\sin(\pi - \alpha) = \sin\alpha$, các hàm khác đổi dấu.\\
        - \textbf{Phụ chéo:} Đổi hàm chéo nhau: $\sin \leftrightarrow \cos, \tan \leftrightarrow \cot$.\\
        - \textbf{Hơn kém pi tang:} $\tan$ và $\cot$ giữ nguyên: $\tan(\alpha + \pi) = \tan\alpha$.
    \end{scaffoldbox}
    \item \textit{Báo cáo, thảo luận:} Học sinh đọc đồng thanh câu vè ghi nhớ, 1 học sinh vận dụng tính nhanh $\sin(\pi - 30^\circ) = \sin 30^\circ = \dfrac{1}{2}$; $\cos(120^\circ) = \cos(180^\circ - 60^\circ) = -\cos 60^\circ = -\dfrac{1}{2}$.
    \item \textit{Kết luận, nhận định:} Giáo viên tổng kết, yêu cầu học sinh đóng khung công thức và áp dụng vào phần luyện tập.
\end{itemize}

\subsubsection*{3. Hoạt động 3: Luyện tập}
\textbf{a) Mục tiêu:} Rèn luyện thành thạo kỹ năng tính các giá trị lượng giác còn lại khi biết một giá trị và khoảng của góc; kỹ năng rút gọn biểu thức lượng giác; kiểm tra kết quả bằng MTCT.

\textbf{b) Nội dung:} Học sinh thực hiện Phiếu học tập số 3 (Worksheet 3):
\begin{itemize}[leftmargin=0.6cm]
    \item \textbf{Bài 1:} Cho $\sin\alpha = \dfrac{3}{5}$ với $\dfrac{\pi}{2} < \alpha < \pi$. Tính $\cos\alpha, \tan\alpha, \cot\alpha$.
    \item \textbf{Bài 2:} Rút gọn biểu thức lượng giác:
    $$A = \sin(\pi - \alpha) + \cos\left(\dfrac{\pi}{2} - \alpha\right) - \tan(\pi + \alpha) \cdot \cot(-\alpha)$$
    \item \textbf{Bài 3:} Chứng minh đẳng thức sau với các góc thỏa mãn điều kiện xác định:
    $$\dfrac{\sin\alpha}{1 + \cos\alpha} + \dfrac{1 + \cos\alpha}{\sin\alpha} = \dfrac{2}{\sin\alpha}$$
\end{itemize}

\textbf{c) Sản phẩm:} Lời giải chi tiết:
\begin{itemize}[leftmargin=0.6cm]
    \item \textbf{Bài 1:}
    \begin{itemize}[leftmargin=0.5cm]
        \item Ta có: $\sin^2\alpha + \cos^2\alpha = 1 \implies \cos^2\alpha = 1 - \sin^2\alpha = 1 - \left(\dfrac{3}{5}\right)^2 = \dfrac{16}{25}$.
        \item Vì $\dfrac{\pi}{2} < \alpha < \pi$ (góc phần tư thứ II) nên $\cos\alpha < 0$. Do đó:
        $$\cos\alpha = -\sqrt{\dfrac{16}{25}} = -\dfrac{4}{5}$$
        \item Tính $\tan\alpha = \dfrac{\sin\alpha}{\cos\alpha} = \dfrac{\dfrac{3}{5}}{-\dfrac{4}{5}} = -\dfrac{3}{4}$.
        \item Tính $\cot\alpha = \dfrac{1}{\tan\alpha} = -\dfrac{4}{3}$.
    \end{itemize}
    \item \textbf{Bài 2:}
    \begin{itemize}[leftmargin=0.5cm]
        \item Áp dụng công thức cung liên kết:
        $\sin(\pi - \alpha) = \sin\alpha$; $\cos\left(\dfrac{\pi}{2} - \alpha\right) = \sin\alpha$;
        $\tan(\pi + \alpha) = \tan\alpha$; $\cot(-\alpha) = -\cot\alpha$.
        \item Thay vào biểu thức $A$:
        $$A = \sin\alpha + \sin\alpha - [\tan\alpha \cdot (-\cot\alpha)] = 2\sin\alpha - (-\tan\alpha \cdot \cot\alpha) = 2\sin\alpha - (-1) = 2\sin\alpha + 1$$
    \end{itemize}
    \item \textbf{Bài 3:} Biến đổi vế trái:
    $$\text{VT} = \dfrac{\sin^2\alpha + (1 + \cos\alpha)^2}{\sin\alpha(1 + \cos\alpha)} = \dfrac{\sin^2\alpha + 1 + 2\cos\alpha + \cos^2\alpha}{\sin\alpha(1 + \cos\alpha)}$$
    Vì $\sin^2\alpha + \cos^2\alpha = 1$ nên tử số là:
    $$(\sin^2\alpha + \cos^2\alpha) + 1 + 2\cos\alpha = 1 + 1 + 2\cos\alpha = 2 + 2\cos\alpha = 2(1 + \cos\alpha)$$
    Do đó:
    $$\text{VT} = \dfrac{2(1 + \cos\alpha)}{\sin\alpha(1 + \cos\alpha)} = \dfrac{2}{\sin\alpha} = \text{VP} \quad (\text{đpcm})$$
\end{itemize}

\textbf{d) Tổ chức thực hiện:}
\begin{itemize}[leftmargin=0.6cm]
    \item \textit{Chuyển giao nhiệm vụ:} Học sinh làm bài theo nhóm đôi, thực hiện tính toán vào phiếu học tập.
    \item \textit{Thực hiện nhiệm vụ (Hỗ trợ HS TB-Yếu):}
    \begin{scaffoldbox}{Giàn giáo sư phạm 3.3: Lỗi sai học sinh cần tránh ở Bài 1}
        \textbf{Lưu ý sống còn:} Từ $\cos^2\alpha = \dfrac{16}{25}$, ta có 2 giá trị $\pm \dfrac{4}{5}$. Học sinh trung bình rất hay quên xét góc phần tư để chọn dấu trừ! Bắt buộc nhìn điều kiện $\dfrac{\pi}{2} < \alpha < \pi$ trước khi khai căn!
    \end{scaffoldbox}
    \item \textit{Báo cáo, thảo luận:} 3 cặp đôi lên bảng trình bày 3 bài toán. Các cặp đôi khác nhận xét, đối chiếu kết quả.
    \item \textit{Kết luận, nhận định:} Giáo viên nhận xét, nhấn mạnh kỹ năng khai căn kết hợp điều kiện dấu và phương pháp quy đồng trong rút gọn lượng giác.
\end{itemize}

\subsubsection*{4. Hoạt động 4: Vận dụng}
\textbf{a) Mục tiêu:} Vận dụng giá trị lượng giác của góc để giải quyết bài toán kỹ thuật tính toán góc nghiêng của tấm pin năng lượng mặt trời.

\textbf{b) Nội dung:}
Để thu được năng lượng mặt trời tối đa, một tấm pin năng lượng mặt trời được lắp đặt nghiêng một góc $\beta$ so với mặt phẳng nằm ngang. Khi mặt trời chiếu góc nghiêng tia sáng $\alpha = 60^\circ$ so với phương ngang, lượng bức xạ $I$ nhận được tỷ lệ thuận với $\sin(\alpha + \beta)$ theo công thức $I = I_0 \cdot \sin(\alpha + \beta)$.
\begin{itemize}[leftmargin=0.6cm]
    \item Nếu tấm pin được đặt nằm ngang ($\beta = 0^\circ$), tính lượng bức xạ $I$ theo $I_0$.
    \item Người ta muốn lượng bức xạ đạt cực đại ($I = I_0$), tức là tia sáng vuông góc với tấm pin. Hỏi cần điều chỉnh góc nghiêng $\beta$ của tấm pin bằng bao nhiêu độ?
\end{itemize}

\textbf{c) Sản phẩm:} Lời giải:
\begin{itemize}[leftmargin=0.6cm]
    \item Khi $\beta = 0^\circ$: $I = I_0 \cdot \sin(60^\circ + 0^\circ) = I_0 \cdot \sin 60^\circ = I_0 \cdot \dfrac{\sqrt{3}}{2} \approx 0{,}866 I_0$.
    \item Để $I = I_0 \iff \sin(\alpha + \beta) = 1 \iff \alpha + \beta = 90^\circ \iff 60^\circ + \beta = 90^\circ \iff \beta = 30^\circ$.
    \item Vậy góc nghiêng của tấm pin cần lắp đặt là $\beta = 30^\circ$.
\end{itemize}

\textbf{d) Tổ chức thực hiện:}
\begin{itemize}[leftmargin=0.6cm]
    \item \textit{Chuyển giao nhiệm vụ:} Giáo viên giao nhiệm vụ liên môn Vật lý -- Kỹ thuật STEM cho cả lớp.
    \item \textit{Thực hiện nhiệm vụ:} Học sinh áp dụng bảng giá trị lượng giác và tính chất $\sin \le 1$.
    \item \textit{Báo cáo, thảo luận:} 1 học sinh trả lời nhanh tại chỗ, giải thích ý nghĩa kinh tế -- thực tiễn của việc điều chỉnh góc pin mặt trời theo mùa.
    \item \textit{Kết luận, nhận định:} Giáo viên tổng kết toàn bộ Bài 1: Nắm chắc góc lượng giác, đường tròn lượng giác, bảng dấu và 4 hệ thức cơ bản. Giao bài tập về nhà trong SGK.
\end{itemize}

% =========================================================================
% PHỤ LỤC
% =========================================================================
\newpage
\section*{IV. PHỤ LỤC HỌC LIỆU BỔ TRỢ}

\subsection*{1. Hệ thống Phiếu học tập (Worksheet) phân tầng bậc thang}

\begin{pedabox}{PHIẾU HỌC TẬP SỐ 1: GÓC VÀ ĐƠN VỊ RADIAN (Phân tầng 3 mức độ)}
\textbf{Họ và tên học sinh:} \dotfill \textbf{Lớp:} 11A\dots \ \textbf{Ngày:} \dotfill

\textbf{MỨC 1: NHẬN BIẾT (Điền khuyết có giàn giáo)}
\begin{enumerate}[leftmargin=0.5cm]
    \item Điền số thích hợp vào chỗ trống:\\
    a) $180^\circ = \dots\dots \text{ rad}$; \quad b) $90^\circ = \dfrac{\pi}{\dots\dots} \text{ rad}$; \quad c) $\dfrac{\pi}{6} \text{ rad} = \dots\dots^\circ$; \quad d) $\dfrac{\pi}{4} \text{ rad} = \dots\dots^\circ$.
    \item Điền dấu ($+$ hoặc $-$): Chiều ngược chiều kim đồng hồ là chiều \dotfill; chiều cùng chiều kim đồng hồ là chiều \dotfill.
\end{enumerate}

\textbf{MỨC 2: THÔNG HIỂU (Áp dụng công thức)}
\begin{enumerate}[leftmargin=0.5cm]
    \item Đổi các số đo góc sau sang radian: $120^\circ; -135^\circ; 210^\circ$.\\
    \textit{Hướng dẫn giàn giáo:} Sử dụng công thức $\alpha = a^\circ \cdot \dfrac{\pi}{180^\circ}$.
    \item Cho đường tròn bán kính $R = 8\text{ cm}$. Tính độ dài cung tròn có số đo $\alpha = \dfrac{3\pi}{4}\text{ rad}$.
\end{enumerate}

\textbf{MỨC 3: VẬN DỤNG (Liên hệ thực tiễn)}
\begin{enumerate}[leftmargin=0.5cm]
    \item Kim giờ của đồng hồ có chiều dài $8\text{ cm}$. Trong thời gian 6 giờ, đầu mút kim giờ quét một cung có độ dài bao nhiêu cm?
\end{enumerate}
\end{pedabox}

\vspace{5pt}

\begin{pedabox}{PHIẾU HỌC TẬP SỐ 2: ĐƯỜNG TRÒN LƯỢNG GIÁC \& GIÁ TRỊ LƯỢNG GIÁC}
\textbf{MỨC 1: NHẬN BIẾT}
\begin{enumerate}[leftmargin=0.5cm]
    \item Trên đường tròn lượng giác, trục hoành $Ox$ gọi là trục \dotfill, trục tung $Oy$ gọi là trục \dotfill.
    \item Xác định dấu ($>$ hoặc $<$) thích hợp:
    $\sin 120^\circ \dots\dots 0$; \quad $\cos 120^\circ \dots\dots 0$; \quad $\tan 240^\circ \dots\dots 0$; \quad $\cos\left(-45^\circ\right) \dots\dots 0$.
\end{enumerate}

\textbf{MỨC 2: THÔNG HIỂU}
\begin{enumerate}[leftmargin=0.5cm]
    \item Sử dụng MTCT Casio fx-580VN X, tính giá trị chính xác hoặc làm tròn 4 chữ số thập phân:
    a) $\sin\dfrac{7\pi}{6}$; \quad b) $\cos\left(-\dfrac{3\pi}{4}\right)$; \quad c) $\tan 195^\circ$; \quad d) $\cot\dfrac{5\pi}{3}$.
    \item Tìm số đo góc $\alpha \in \left[\dfrac{\pi}{2}; \pi\right]$ biết $\sin\alpha = \dfrac{1}{2}$.
\end{enumerate}

\textbf{MỨC 3: VẬN DỤNG}
\begin{enumerate}[leftmargin=0.5cm]
    \item Cho góc $\alpha$ thỏa mãn $\tan\alpha = -2$ và $\dfrac{3\pi}{2} < \alpha < 2\pi$. Không dùng máy tính cầm tay, hãy xác định dấu của biểu thức $P = \sin\alpha \cdot \cos\alpha + \cos^2\alpha$.
\end{enumerate}
\end{pedabox}

\vspace{5pt}

\begin{pedabox}{PHIẾU HỌC TẬP SỐ 3: HỆ THỨC CƠ BẢN \& CUNG LIÊN KẾT}
\textbf{MỨC 1: NHẬN BIẾT}
\begin{enumerate}[leftmargin=0.5cm]
    \item Viết lại 4 hệ thức lượng giác cơ bản:
    a) $\sin^2\alpha + \cos^2\alpha = \dots\dots$; \quad b) $1 + \tan^2\alpha = \dots\dots$; \quad c) $1 + \cot^2\alpha = \dots\dots$; \quad d) $\tan\alpha \cdot \cot\alpha = \dots\dots$
    \item Điền vào chỗ trống: $\sin(\pi - \alpha) = \dots\dots$; \quad $\cos(-\alpha) = \dots\dots$; \quad $\tan(\pi + \alpha) = \dots\dots$.
\end{enumerate}

\textbf{MỨC 2: THÔNG HIỂU}
\begin{enumerate}[leftmargin=0.5cm]
    \item Cho $\cos\alpha = -\dfrac{5}{13}$ với $\pi < \alpha < \dfrac{3\pi}{2}$. Tính $\sin\alpha, \tan\alpha, \cot\alpha$.\\
    \textit{Giàn giáo gợi ý:} Bước 1: Tính $\sin^2\alpha = 1 - \cos^2\alpha$. Bước 2: Vì $\pi < \alpha < \dfrac{3\pi}{2} \implies \sin\alpha < 0 \implies \sin\alpha = -\sqrt{\dots}$.
    \item Rút gọn biểu thức: $M = \cos(\pi - x) + \sin\left(x + \dfrac{\pi}{2}\right) - \tan\left(\dfrac{\pi}{2} - x\right) \cdot \tan x$.
\end{enumerate}

\textbf{MỨC 3: VẬN DỤNG}
\begin{enumerate}[leftmargin=0.5cm]
    \item Chứng minh rằng với mọi góc $\alpha$ làm cho biểu thức xác định, ta có:
    $$\dfrac{1 - 2\sin\alpha\cos\alpha}{\cos^2\alpha - \sin^2\alpha} = \dfrac{\cos\alpha - \sin\alpha}{\cos\alpha + \sin\alpha}$$
\end{enumerate}
\end{pedabox}

\subsection*{2. Rubric đánh giá năng lực học sinh theo 4 mức độ}

\begin{table}[h!]
\centering
\small
\begin{tabularx}{\textwidth}{|p{2.8cm}|X|X|X|X|}
\hline
\textbf{Tiêu chí} & \textbf{Chưa đạt (Mức 1)} & \textbf{Đạt (Mức 2)} & \textbf{Khá (Mức 3)} & \textbf{Tốt (Mức 4)} \\
\hline
\textbf{1. Khái niệm góc và Đơn vị đo} & Chưa đổi được độ sang radian, nhầm lẫn chiều quay lượng giác. & Đổi được góc đặc biệt giữa độ và radian; vẽ được góc hình học cơ bản. & Đổi thành thạo góc bất kỳ; áp dụng đúng công thức tính độ dài cung tròn $l = \alpha R$. & Giải thích bản chất radian, vận dụng linh hoạt giải bài toán chuyển động tròn thực tế. \\
\hline
\textbf{2. Đường tròn \& Giá trị lượng giác} & Chưa xác định được điểm biểu diễn trên đường tròn đơn vị, nhầm trục sin và cos. & Xác định được điểm biểu diễn góc đặc biệt; tra cứu được bảng dấu cơ bản. & Xác định chính xác dấu 4 hàm số trên 4 góc phần tư; tính đúng giá trị lượng giác góc đặc biệt. & Thành thạo lập luận hình học để tìm giá trị lượng giác; giải thích sâu sắc ý nghĩa tọa độ. \\
\hline
\textbf{3. Hệ thức cơ bản \& Cung liên kết} & Chưa thuộc công thức cơ bản; sai lầm khi chọn dấu khai căn $\cos^2 = 1 - \sin^2$. & Thuộc công thức; tính được các giá trị còn lại khi góc ở góc phần tư thứ nhất. & Tính đúng các giá trị còn lại khi góc ở góc phần tư II, III, IV; rút gọn biểu thức cơ bản. & Chứng minh được các đẳng thức lượng giác nâng cao; giải quyết trọn vẹn bài toán thực tiễn. \\
\hline
\textbf{4. Năng lực số \& Năng lực AI} & Chưa biết chỉnh máy tính sang chế độ Radian; chưa biết đặt prompt cho AI. & Thao tác được phím đổi đơn vị trên MTCT; đặt được câu hỏi cơ bản cho AI. & Khai thác thành thạo MTCT kiểm tra kết quả; viết prompt yêu cầu AI giải thích bảng dấu. & Làm chủ hoàn toàn MTCT; phân tích, phát hiện được sai sót hoặc giới hạn của AI trong bài toán lượng giác. \\
\hline
\end{tabularx}
\end{table}

\end{document}